\documentclass[aspectratio=169, 10pt]{beamer}
\usepackage{sunass_estilo}

\title{Módulo VII: Competencia Imperfecta}
\subtitle{Guía Teórica de Repaso Rápido · PEU SUNASS}
\author[PEU SUNASS · Economía]{Preparación Examen de Admisión PEU SUNASS}
\institute{Basado en el Temario Oficial de Economía · Organización Industrial}
\date[2026]{2026}

\begin{document}

\begin{frame}[plain]
  \titlepage
\end{frame}

% ==============================================================================
% SLIDE 2: MAPA CONCEPTUAL
% ==============================================================================
\begin{frame}{Mapa Conceptual: Competencia Imperfecta}
  \vspace{0.15cm}
  \begin{columns}[T]
    \begin{column}{0.32\textwidth}
      \tarjeta{azulPrincipal}{1. Teoría de Juegos}{
        Interdependencia estratégica:\\
        • Juegos estáticos (Nash, estrategias dominantes, Dilema del Prisionero).\\
        • Juegos dinámicos (inducción hacia atrás, ENPS) y repetidos.
      }
    \end{column}
    \begin{column}{0.32\textwidth}
      \tarjeta{naranjaAcento}{2. Oligopolio}{
        • \textbf{Cournot:} cantidades simultáneas.\\
        • \textbf{Stackelberg:} líder-seguidor en \(q\).\\
        • \textbf{Bertrand:} precios simultáneos.\\
        • \textbf{Cártel / Colusión:} inestabilidad.
      }
    \end{column}
    \begin{column}{0.32\textwidth}
      \tarjeta{verdeAcento}{3. Competencia Monopolística}{
        Muchas firmas con productos \textbf{diferenciados} y libre entrada.\\
        • Corto plazo: \(\pi > 0\).\\
        • Largo plazo: \(\pi = 0\) con \textbf{exceso de capacidad}.
      }
    \end{column}
  \end{columns}

  \vspace{0.2cm}
  \recuadroclave{
    \small\textbf{Ordenamiento de Cantidades y Precios (Bienes Homogéneos):} \quad \(Q_{\text{Cártel}} < Q_{\text{Cournot}} < Q_{\text{Stackelberg}} < Q_{\text{Bertrand}} = Q_{\text{Competencia}}\)
  }
  \mensajefinal{En oligopolio, el bienestar social crece a medida que las empresas compiten más agresivamente (de colusión hacia Bertrand).}
\end{frame}

% ==============================================================================
% SLIDE 3: TEORÍA DE JUEGOS (NASH, DILEMA, DINÁMICOS Y REPETIDOS)
% ==============================================================================
\begin{frame}{1. Teoría de Juegos: Nash, Dinámicos y Repetidos}
  \vspace{0.02cm}
  \begin{columns}[T]
    \begin{column}{0.32\textwidth}
      \tarjeta{azulPrincipal}{Equilibrio de Nash (EN)}{
        Perfil de estrategias \((s_1^*, s_2^*)\) donde cada jugador elige su \textbf{mejor respuesta} dada la estrategia del rival: nadie tiene incentivo a desviarse unilateralmente.\\[2pt]
        • \textbf{Estrategia dominante:} es óptima sin importar qué haga el rival.
      }
    \end{column}
    \begin{column}{0.32\textwidth}
      \tarjeta{rojoAcento}{Dilema del Prisionero y ENPS}{
        • \textbf{Dilema del Prisionero:} El EN en estrategias dominantes es Pareto-inferior a la cooperación (explica la inestabilidad de los cárteles).\\
        • \textbf{Juegos Secuenciales:} Se resuelven por \textbf{inducción hacia atrás} para eliminar amenazas no creíbles (\textbf{ENPS}).
      }
    \end{column}
    \begin{column}{0.32\textwidth}
      \tarjeta{verdeAcento}{Juegos Repetidos (Folk Theorem)}{
        Con horizonte finito \(T\) conocido, por inducción hacia atrás no hay cooperación.\\[2pt]
        Con \textbf{horizonte infinito} (o probabilidad de continuación \(\delta\)), se sostiene la colusión con estrategia \textit{Grim Trigger} si \(\delta \ge \delta^*\).
      }
    \end{column}
  \end{columns}

  \vspace{0.05cm}
  \recuadroclave{
    \small\textbf{Estrategias Mixtas:} Si no existe EN en estrategias puras, todo juego finito tiene al menos un EN en estrategias mixtas (haciendo indiferente al rival entre sus acciones puras).
  }
  \mensajefinal{Todo equilibrio en estrategias dominantes es de Nash, pero no todo Equilibrio de Nash usa estrategias dominantes.}
\end{frame}

% ==============================================================================
% SLIDE 4: MODELO DE COURNOT Y STACKELBERG
% ==============================================================================
\begin{frame}{2. Oligopolio en Cantidades: Cournot y Stackelberg}
  \begin{columns}[T]
    \begin{column}{0.52\textwidth}
      {\small\textbf{Funciones de Reacción en Duopolio de Cournot}}
      \vspace{0.05cm}
      \begin{tikzpicture}[scale=0.78]
        \begin{axis}[
          axis lines=left,
          xlabel={$q_1$}, ylabel={$q_2$},
          xmin=0, xmax=9.8, ymin=0, ymax=9.8,
          xtick={3.0, 4.5, 9.0}, xticklabels={$q_1^C$, $q_1^S\!=\!q_1^{mon}$, $Q_C$},
          ytick={2.15, 3.05, 4.5, 9.0}, yticklabels={$q_2^S$, $q_2^C$, $q_2^{mon}$, $Q_C$},
          tick label style={font=\tiny},
          width=8.0cm, height=5.2cm,
          clip=false
        ]
          % R1: q1 = 4.5 - 0.5 q2 => q2 = 9 - 2 q1 (corte en q1=4.5, q2=0; q1=0, q2=9)
          \addplot[azulPrincipal, very thick, domain=0:4.5] {9 - 2*x};
          \node[azulPrincipal, font=\scriptsize, anchor=west] at (axis cs:0.4, 8.4) {$R_1(q_2)$};
          % R2: q2 = 4.5 - 0.5 q1 (corte en q1=9, q2=0; q1=0, q2=4.5)
          \addplot[rojoAcento, very thick, domain=0:9.0] {4.5 - 0.5*x};
          \node[rojoAcento, font=\scriptsize, anchor=south] at (axis cs:7.4, 1.0) {$R_2(q_1)$};
          % Cournot (3,3)
          \addplot[only marks, mark=*, mark size=2.8pt, verdeAcento] coordinates {(3, 3)};
          \node[verdeAcento, font=\scriptsize, anchor=south west] at (axis cs:3.05, 3.1) {\textbf{Cournot}};
          % Stackelberg (4.5, 2.25)
          \addplot[only marks, mark=square*, mark size=2.6pt, naranjaAcento] coordinates {(4.5, 2.25)};
          \node[naranjaAcento, font=\scriptsize, anchor=south west] at (axis cs:4.55, 2.35) {\textbf{Stackelberg}};
          \draw[dotted, grisTexto] (axis cs:3,0) -- (axis cs:3,3) -- (axis cs:0,3);
          \draw[dotted, grisTexto] (axis cs:4.5,0) -- (axis cs:4.5,2.25) -- (axis cs:0,2.25);
        \end{axis}
      \end{tikzpicture}
    \end{column}

    \begin{column}{0.45\textwidth}
      {\scriptsize\color{grisTexto}Cournot (Simultáneo)}\\[1pt]
      {\normalsize\color{verdeAcento}\textbf{\(q_1^* = q_2^* = \frac{a-c}{3b} \implies Q = \frac{2}{3}Q_C\)}}\\[4pt]
      {\scriptsize
        • Demanda \(P = a - b(q_1 + q_2)\) y \(CMg = c\).\\
        • Cada firma elige \(q_i\) tomando \(q_j\) como dado (\textbf{sustitutos estratégicos}: pendientes negativas).\\
        • Con \(N\) firmas: \(Q_N = \frac{N}{N+1}Q_C\) (converge a competencia cuando \(N \to \infty\)).\\[4pt]
      }
      \indicador{Stackelberg (Secuencial)}{naranjaAcento}{\(q_1^S = \frac{a-c}{2b}, \quad q_2^S = \frac{a-c}{4b}\)}
      {\scriptsize
        • La firma 1 (líder) incorpora \(R_2(q_1)\) en su beneficio.\\
        • \textbf{Ventaja del primer movedor:} \(q_1^S > q_1^C\) y \(\pi_1^S > \pi_1^C\).
      }
    \end{column}
  \end{columns}

  \mensajefinal{En Cournot con \(N\) firmas, el Índice de Lerner promedio es \(L = \frac{HHI}{|\varepsilon_d|}\), vinculando concentración (\(HHI\)) y poder de mercado.}
\end{frame}

% ==============================================================================
% SLIDE 5: OLIGOPOLIO ASIMÉTRICO, CÁRTEL Y COMPARATIVA DE EFICIENCIA (PES)
% ==============================================================================
\begin{frame}{2b. Oligopolio Asimétrico, Cártel y Pérdida de Eficiencia}
  \begin{columns}[T]
    \begin{column}{0.51\textwidth}
      {\small\textbf{Jerarquía de Bienestar (\(P = a - bQ\), \(CMg = c\))}}
      \vspace{0.05cm}
      \begin{tikzpicture}[scale=0.77]
        \begin{axis}[
          axis lines=left,
          xlabel={Cantidad ($Q$)}, ylabel={Precio ($P$)},
          xmin=0, xmax=9.2, ymin=0, ymax=8.8,
          xtick={3.0, 4.0, 4.5, 6.0}, xticklabels={$Q_{\text{Col}}$, $Q_{\text{Co}}$, $Q_{\text{St}}$, $Q_{\text{Be}}$},
          ytick={2.0, 3.5, 4.0, 5.0, 8.0}, yticklabels={$c$, $P_{\text{Sta}}$, $P_{\text{Cou}}$, $P_{\text{Col}}$, $a$},
          tick label style={font=\tiny},
          width=8.0cm, height=5.1cm, clip=false
        ]
          % Demanda P = 8 - Q, CMg = 2 => QC = 6
          % Q_Col = 3, P_Col = 5; Q_Cou = 4, P_Cou = 4; Q_Sta = 4.5, P_Sta = 3.5; Q_Ber = 6, P_Ber = 2
          \addplot[fill=rojoAcento!16, draw=none] coordinates {(3.0,2.0) (6.0,2.0) (3.0,5.0)} \closedcycle;
          \addplot[fill=naranjaAcento!20, draw=none] coordinates {(4.0,2.0) (6.0,2.0) (4.0,4.0)} \closedcycle;
          \addplot[fill=verdeAcento!24, draw=none] coordinates {(4.5,2.0) (6.0,2.0) (4.5,3.5)} \closedcycle;

          \addplot[azulPrincipal, very thick, domain=0:7.6] {8 - x};
          \node[azulPrincipal, font=\scriptsize, anchor=west] at (axis cs:7.0, 1.3) {$D$};
          \addplot[rojoAcento, very thick, domain=0:8.3] {2};
          \node[rojoAcento, font=\scriptsize, anchor=west] at (axis cs:8.35, 2.0) {$CMg=c$};

          \draw[dotted, grisTexto] (axis cs:3.0, 0) -- (axis cs:3.0, 5.0) -- (axis cs:0, 5.0);
          \draw[dotted, grisTexto] (axis cs:4.0, 0) -- (axis cs:4.0, 4.0) -- (axis cs:0, 4.0);
          \draw[dotted, grisTexto] (axis cs:4.5, 0) -- (axis cs:4.5, 3.5) -- (axis cs:0, 3.5);
          \addplot[only marks, mark=*, mark size=2pt, azulOscuro] coordinates {(3,5) (4,4) (4.5,3.5) (6,2)};

          \node[rojoAcento, font=\tiny\bfseries, anchor=west] at (axis cs:3.1, 5.35) {Colusión \(\big(\frac{(a-c)^2}{8b}\big)\)};
          \node[naranjaAcento!90!black, font=\tiny\bfseries, anchor=west] at (axis cs:4.1, 4.25) {Cournot \(\big(\frac{(a-c)^2}{18b}\big)\)};
          \node[verdeAcento!90!black, font=\tiny\bfseries, anchor=west] at (axis cs:4.65, 3.55) {Stackelberg \(\big(\frac{(a-c)^2}{32b}\big)\)};
          \node[azulPrincipal, font=\tiny\bfseries, anchor=south west] at (axis cs:6.05, 2.2) {Bertrand (\(PES=0\))};
        \end{axis}
      \end{tikzpicture}
    \end{column}

    \begin{column}{0.46\textwidth}
      {\scriptsize\color{grisTexto}Cournot con Costos Asimétricos (\(c_1 < c_2\))}\\[1pt]
      {\small\color{azulPrincipal}\textbf{\(q_1^* = \frac{a-2c_1+c_2}{3b},\quad q_2^* = \frac{a-2c_2+c_1}{3b}\)}}\\[1pt]
      {\scriptsize
        • Precio de equilibrio: \(P^* = \frac{a + c_1 + c_2}{3}\) (menor \(CMg \Rightarrow\) mayor cuota).\\[3pt]
      }
      {\scriptsize\color{grisTexto}Stackelberg Asimétrico (Líder: menor \(CMg\) o mayor \(CF\))}\\[1pt]
      {\small\color{naranjaAcento}\textbf{\(q_1^S = \frac{a-2c_1+c_2}{2b},\quad q_2^S = \frac{a-3c_2+2c_1}{4b}\)}}\\[1pt]
      {\scriptsize
        • Precio: \(P^S = \frac{a + 2c_1 + c_2}{4}\) (solución interior si \(a+2c_1 > 3c_2\)).\\[3pt]
      }
      {\scriptsize\color{grisTexto}Cártel con Cuotas (\(q_1 = XQ,\; q_2 = YQ,\; X+Y=1\))}\\[1pt]
      {\small\color{rojoAcento}\textbf{\(Q_{\text{Col}} = \frac{a - Xc_1 - Yc_2}{2b},\quad P_{\text{Col}} = \frac{a + Xc_1 + Yc_2}{2}\)}}\\[1pt]
      {\scriptsize
        • Sin pacto de cuotas y \(c_1 < c_2\), el cártel cierra la planta ineficiente y produce solo con la firma \(1\).
      }
    \end{column}
  \end{columns}

  \mensajefinal{A mayor competencia (\text{Colusión} \(\to\) \text{Cournot} \(\to\) \text{Stackelberg} \(\to\) \text{Bertrand}), menor \(PES\).}
\end{frame}

% ==============================================================================
% SLIDE 6: MODELO DE BERTRAND HOMOGÉNEO Y COLUSIÓN (CÁRTEL)
% ==============================================================================
\begin{frame}{3. Modelo de Bertrand (Homogéneo) y Colusión}
  \begin{columns}[T]
    \begin{column}{0.49\textwidth}
      {\small\textbf{Funciones de Reacción en Precios (\textit{Undercutting})}}
      \vspace{0.03cm}
      \begin{tikzpicture}[scale=0.68, >=stealth]
        % Ejes
        \draw[->, thick] (0,0) -- (6.2,0) node[right, font=\scriptsize] {$p_1$};
        \draw[->, thick] (0,0) -- (0,5.8) node[above, font=\scriptsize] {$p_2$};
        % Diagonal 45°
        \draw[dashed, thin, grisTexto] (0,0) -- (5.3,5.3);
        % Guías p_max y CMg
        \draw[dotted, grisTexto] (4.3,0) node[below, font=\tiny, black] {$p_1^{\max}$} -- (4.3,4.55);
        \draw[dotted, grisTexto] (0,4.3) node[left, font=\tiny, black] {$p_2^{\max}$} -- (4.55,4.3);
        % FR1 (azul)
        \draw[azulPrincipal, very thick] (1.4,0) -- (1.4,1.65) -- (4.3,4.55) -- (4.3,5.4) node[right, font=\scriptsize] {$FR_1$};
        % FR2 (verde)
        \draw[verdeAcento, very thick] (0,1.4) -- (1.65,1.4) -- (4.55,4.3) -- (5.5,4.3) node[above, font=\scriptsize] {$FR_2$};
        % Nash equilibrium at (1.4, 1.4)
        \fill[rojoAcento] (1.4,1.4) circle (2.8pt);
        \node[below, font=\tiny, black] at (1.4,0) {$CMg_1$};
        \node[left, font=\tiny, black] at (0,1.4) {$CMg_2$};
        \node[rojoAcento, font=\tiny\bfseries, anchor=west] at (1.65,0.8) {EN: \((CMg, CMg)\)};
      \end{tikzpicture}\\[-2pt]
      {\tiny\textbf{Demanda discontinua:} \(D_i = D(p_i)\) si \(p_i < p_j\);\; \(\frac{1}{2}D(p_i)\) si \(p_i = p_j\);\; \(0\) si \(p_i > p_j\).}
    \end{column}

    \begin{column}{0.48\textwidth}
      \tarjeta{azulPrincipal}{Paradoja de Bertrand (Bienes Homogéneos)}{
        Con bienes idénticos, \(CMg = c\) cte. y sin límites de capacidad: el único Equilibrio de Nash es \(\mathbf{P_1^* = P_2^* = c}\) (\(\pi = 0\)).\\[2pt]
        • \textbf{Mecanismo:} mientras \(p_j > c\), conviene recortar un céntimo (\(p_i = p_j - \varepsilon\)) y captar todo el mercado.\\
        • \textbf{Cómo se rompe:} diferenciación, restricción de capacidad (Edgeworth) o repetición infinita.
      }
      \vspace{0.04cm}
      \tarjeta{rojoAcento}{Colusión Tácita y Estabilidad del Cártel}{
        • \textbf{Grim Trigger (Bertrand homogéneo con \(N\) firmas):} la colusión se sostiene si \(\delta \ge 1 - \frac{1}{N}\).\\
        • Si \(P_M > CMg\), cada firma tiene incentivo a desviarse (\textit{Dilema del Prisionero}).
      }
    \end{column}
  \end{columns}

  \mensajefinal{En Bertrand homogéneo bastan dos empresas para replicar el precio competitivo (\(P^* = c\)).}
\end{frame}

% ==============================================================================
% SLIDE 7: OLIGOPOLIO CON PRODUCTOS DIFERENCIADOS (COURNOT Y BERTRAND)
% ==============================================================================
\begin{frame}{3b. Oligopolio Diferenciado: Cournot y Bertrand}
  \begin{columns}[T]
    \begin{column}{0.49\textwidth}
      {\small\textbf{Bertrand Diferenciado (Complementos Estratégicos)}}
      \vspace{0.05cm}
      \begin{tikzpicture}[scale=0.76]
        \begin{axis}[
          axis lines=left,
          xlabel={$p_1$}, ylabel={$p_2$},
          xmin=0, xmax=12.5, ymin=0, ymax=12.5,
          xtick={6, 8}, xticklabels={$p_1^0$, $p_1^*$},
          ytick={6, 8}, yticklabels={$p_2^0$, $p_2^*$},
          width=7.8cm, height=5.0cm, clip=false
        ]
          % FR1: p1 = 6 + 0.25*p2 => (6,0) a (9.1, 12.4)
          \addplot[azulPrincipal, very thick, domain=6:9.0] {4*(x - 6)};
          \node[azulPrincipal, font=\scriptsize, anchor=south] at (axis cs:9.0, 12.0) {$FR_1(p_2)$};
          % FR2: p2 = 6 + 0.25*p1 => (0,6) a (12,9)
          \addplot[rojoAcento, very thick, domain=0:12] {6 + 0.25*x};
          \node[rojoAcento, font=\scriptsize, anchor=west] at (axis cs:12.0, 9.0) {$FR_2(p_1)$};
          % Diagonal 45 y equilibrio (8,8)
          \addplot[dashed, grisTexto!65, domain=0:11] {x};
          \draw[dotted, grisTexto] (axis cs:8,0) -- (axis cs:8,8) -- (axis cs:0,8);
          \addplot[only marks, mark=*, mark size=2.6pt, verdeAcento] coordinates {(8,8)};
          \node[verdeAcento, font=\scriptsize\bfseries, anchor=north west] at (axis cs:8.2, 7.15) {EN \((p_1^*, p_2^*)\)};
        \end{axis}
      \end{tikzpicture}
    \end{column}

    \begin{column}{0.48\textwidth}
      {\scriptsize\color{grisTexto}Bertrand Diferenciado (\(q_i = \alpha - \beta p_i + \delta p_j\))}\\[1pt]
      {\small\color{azulPrincipal}\textbf{\(FR_1:\; p_1 = \frac{\alpha + \beta c_1}{2\beta} + \frac{\delta}{2\beta}p_2 \;\; \text{(pendiente} > 0\text{)}\)}}\\[2pt]
      {\scriptsize
        • Si \(c_1 = c_2 = c\): \(\mathbf{p^* = \frac{\alpha + \beta c}{2\beta - \delta} > c}\) (margen positivo).\\
        • Son \textbf{complementos estratégicos}: si el rival sube \(p_j\), conviene subir \(p_i\).\\[4pt]
      }
      {\scriptsize\color{grisTexto}Cournot Diferenciado (\(p_i = a - q_i - d q_j,\; d \in [0, 1]\))}\\[1pt]
      {\small\color{verdeAcento}\textbf{\(q_i^* = \frac{a-c}{2+d},\quad p_i^* = \frac{a+c+cd}{2+d},\quad \pi_i^* = \left(\frac{a-c}{2+d}\right)^2 - CF\)}}\\[2pt]
      {\scriptsize
        • \(d = 0\): monopolios independientes (\(q_i = \frac{a-c}{2}\)).\\
        • \(d = 1\): bienes homogéneos (Cournot clásico).\\
        • Mayor diferenciación (\(\downarrow d\)) \(\implies \uparrow P\) y \(\uparrow \pi\).
      }
    \end{column}
  \end{columns}

  \mensajefinal{La diferenciación de productos (\(d < 1\)) suaviza la competencia en precios (\(\text{FR}\) con pendiente positiva) y rescata márgenes.}
\end{frame}

% ==============================================================================
% SLIDE 6: COMPETENCIA MONOPOLÍSTICA Y EXCESO DE CAPACIDAD
% ==============================================================================
\begin{frame}{4. Competencia Monopolística (Chamberlin)}
  \begin{columns}[T]
    \begin{column}{0.53\textwidth}
      {\small\textbf{Equilibrio a Largo Plazo: Tangencia y Exceso de Capacidad}}
      \vspace{0.05cm}
      \begin{tikzpicture}[scale=0.78]
        \begin{axis}[
          axis lines=left,
          xlabel={Cantidad ($q_i$)}, ylabel={Precio / Costo},
          xmin=0, xmax=9, ymin=0, ymax=8.2,
          xtick={3.5, 5.5}, xticklabels={$q_{LP}^*$, $q_{EME}$},
          ytick={3.8}, yticklabels={$P_{LP}^*$},
          width=8.2cm, height=5.2cm
        ]
          % CMe = 3 + 0.2*(x-5.5)^2 -> en x=3.5: CMe = 3.8, derivada = -0.8
          \addplot[azulPrincipal, very thick, domain=1.0:8.5, samples=50] {3 + 0.2*(x-5.5)^2};
          \node[azulPrincipal, font=\scriptsize, anchor=west] at (axis cs:7.6, 4.4) {$CMe$};
          % CMg exacto = CMe + x*CMe' = 3 + 0.2*(x-5.5)^2 + 0.4*x*(x-5.5) -> pasa por (3.5, 1.0) y (5.5, 3.0)
          \addplot[rojoAcento, very thick, domain=2.6:6.8, samples=40] {3 + 0.2*(x-5.5)^2 + 0.4*x*(x-5.5)};
          \node[rojoAcento, font=\scriptsize, anchor=east] at (axis cs:6.55, 7.2) {$CMg$};
          % Demanda tangente en (3.5, 3.8) con pendiente -0.8: P = 6.6 - 0.8x
          \addplot[verdeAcento, very thick, domain=0.2:8.0] {6.6 - 0.8*x};
          \node[verdeAcento, font=\scriptsize, anchor=south west] at (axis cs:6.7, 1.35) {$d_i$};
          % IMg = 6.6 - 1.6x -> en x=3.5 vale 1.0 = CMg(3.5) > 0
          \addplot[naranjaAcento, thick, dashed, domain=0.2:4.1] {6.6 - 1.6*x};
          \node[naranjaAcento, font=\tiny, anchor=north east] at (axis cs:2.4, 2.75) {$IMg$};
          % Puntos exactos: tangencia (3.5, 3.8) y corte IMg=CMg en (3.5, 1.0)
          \addplot[only marks, mark=*, mark size=2.4pt, rojoAcento] coordinates {(3.5, 3.8) (3.5, 1.0)};
          \draw[dotted, grisTexto] (axis cs:3.5, 0) -- (axis cs:3.5, 3.8) -- (axis cs:0, 3.8);
          \draw[dotted, grisTexto] (axis cs:5.5, 0) -- (axis cs:5.5, 3.0);
          % Exceso capacidad
          \draw[<->, thick, rojoAcento] (axis cs:3.5, 0.45) -- (axis cs:5.5, 0.45) node[midway, above, font=\tiny] {Exceso cap.};
        \end{axis}
      \end{tikzpicture}
    \end{column}

    \begin{column}{0.44\textwidth}
      {\scriptsize\color{grisTexto}Condición de Equilibrio LP}\\[1pt]
      {\normalsize\color{azulPrincipal}\textbf{\(IMg = CMg \quad \text{y} \quad P = CMe \ (\pi = 0)\)}}\\[4pt]
      {\scriptsize
        • La entrada libre desplaza \(d_i\) hacia la izquierda hasta hacerse \textbf{tangente} al \(CMe\) en su tramo decreciente.\\[4pt]
      }
      \indicador{Dos Ineficiencias Típicas}{rojoAcento}{\(P > CMg \quad \text{y} \quad q_{LP}^* < q_{EME}\)}
      {\scriptsize
        • \textbf{1. Margen sobre CMg:} \(P > CMg\) genera poder de mercado.\\
        • \textbf{2. Exceso de capacidad:} produce por debajo de la Escala Mínima Eficiente (\(q_{EME}\)).\\
        • \textit{Compensación social:} mayor \textbf{variedad} de productos para los consumidores.
      }
    \end{column}
  \end{columns}

  \mensajefinal{En competencia monopolística a largo plazo el beneficio económico es cero (\(\pi=0\)), pero NO se produce al costo medio mínimo.}
\end{frame}

% ==============================================================================
% SLIDE 7: EMPRESA DOMINANTE, HOTELLING, MERCADOS DESAFIABLES Y FUSIÓN
% ==============================================================================
\begin{frame}{5. Empresa Dominante, Hotelling y Mercados Desafiables}
  \vspace{0.05cm}
  \begin{columns}[T]
    \begin{column}{0.32\textwidth}
      \tarjeta{azulPrincipal}{Empresa Dominante (Forchheimer)}{
        Líder fija precio \(P\) frente a una franja competitiva tomadora de precios:\\[2pt]
        • Demanda residual: \(D_{\text{dom}}(P) = D(P) - S_{\text{franja}}(P)\).\\
        • El líder iguala \(IMg_{\text{dom}} = CMg_{\text{dom}}\).
      }
    \end{column}
    \begin{column}{0.32\textwidth}
      \tarjeta{naranjaAcento}{Diferenciación Espacial (Hotelling)}{
        • Con \textbf{precios fijos} (exógenos), los rivales se aglomeran en el centro (\textbf{diferenciación mínima}).\\
        • Con costos de transporte \textbf{cuadráticos} y precios libres, maximizan distancia (\textbf{diferenciación máxima}) para suavizar la guerra de precios.
      }
    \end{column}
    \begin{column}{0.32\textwidth}
      \tarjeta{verdeAcento}{Mercados Desafiables (Baumol)}{
        Si hay libre entrada y \textbf{cero costos hundidos} (\textit{hit-and-run}):\\[2pt]
        Incluso un monopolio natural debe fijar \(\mathbf{P = CMe}\) (\(\pi = 0\)) para disuadir la entrada potencial.
      }
    \end{column}
  \end{columns}

  \vspace{0.1cm}
  \recuadroclave{
    \small\textbf{Paradoja de Fusión en Cournot (Salant-Switzer-Reynolds):} Sin ganancias de eficiencia, fusionar 2 firmas en un mercado Cournot con \(N \ge 3\) y \(CMg\) cte. suele \textbf{reducir} sus beneficios conjuntos porque los rivales expanden producción.
  }
  \mensajefinal{No es el número de empresas actuales lo que disciplina a un monopolio sin costos hundidos, sino la amenaza creíble de entrada.}
\end{frame}

% ==============================================================================
% SLIDE 8: CONCLUSIONES Y CHEAT-SHEET MÓDULO VII
% ==============================================================================
\begin{frame}{Conclusiones y Formulario Rápido — Módulo VII}
  \vspace{0.05cm}
  \recuadroclave{
    \small\textbf{Regla de Oro:} La competencia estratégica en cantidades (Cournot) deja márgenes positivos (\(P > CMg\)), mientras en precios con bien homogéneo (Bertrand) colapsa a \(P = CMg\).
  }

  \vspace{0.15cm}
  \begin{columns}[T]
    \begin{column}{0.32\textwidth}
      \tarjeta{verdeAcento}{Fórmulas Imprescindibles}{
        • Cournot \(N\): \(Q = \frac{N}{N+1}Q_C\)\\
        • Stackelberg: \(q_1^S = \frac{1}{2}Q_C,\ q_2^S = \frac{1}{4}Q_C\)\\
        • Bertrand dif.: \(p^* = \frac{\alpha+\beta c}{2\beta-\delta}\)\\
        • Colusión \(N\): \(\delta \ge 1 - \frac{1}{N}\)
      }
    \end{column}
    \begin{column}{0.32\textwidth}
      \tarjeta{azulPrincipal}{Trampas de Examen}{
        • Confundir Cournot con Stackelberg (en Stackelberg solo 1 firma es líder).\\
        • Creer que \(\pi=0\) en comp. monopolística implica \(P=CMg\) (es \(P=CMe > CMg\)).\\
        • Olvidar eliminar amenazas no creíbles (ENPS).
      }
    \end{column}
    \begin{column}{0.32\textwidth}
      \tarjeta{rojoAcento}{Siguientes Pasos}{
        • Practicar Bloque B del Módulo 07 en el Simulador.\\
        • Pasar al Módulo VIII (Externalidades, Bienes Públicos y Asimetría de Información).
      }
    \end{column}
  \end{columns}

  \vspace{0.25cm}
  \begin{center}
    {\color{azulPrincipal}\Large\textbf{¡Éxito!}} \quad {\color{grisTexto}Módulo VII · Competencia Imperfecta · PEU SUNASS 2026}
  \end{center}
\end{frame}

\end{document}
